\section{Introduction}

The comma sequence \oeis{A121805} starts at $1$, and each next term $k'$ of a term $k$ is the
smallest number such that $k' - k$ equals the two-digit number formed by the last digit of $k$ and
the first digit of $k'$. It famously dies after $2137453$ terms at $99999945$. Angelini et
al.~\cite{comma} study this sequence, its variants with other starting values and bases, and the
\emph{child graph}, in which a number may continue with any valid next term rather than the
smallest. They conjecture that all comma sequences in bases $b \ge 3$ are finite and prove it for
$b = 3$.

Dougherty-Bliss and Ter-Saakov~\cite{dbts} proved this conjecture for all bases $3 \le b \le 633$.
Their method reduces the behaviour of comma sequences between consecutive ``danger intervals''
$[d\, b^k - b^2, d\, b^k)$ to a map that depends on $k$ only modulo a period $L(b)$, which yields a
finite graph $G'_b$ that has a cycle if and only if some comma sequence is infinite; they check the
absence of cycles by computer~\cite{dbtscode}. They also prove Conjecture~7.1 of~\cite{comma}.

The child graph has at least one infinite path~\cite[Theorem~5.6]{comma}; the lexicographically
earliest one in base 10 is \oeis{A367620}, starting at the root $20$. Angelini et al.\ determined
its first $30$ branch choices, recorded as \oeis{A399179}, and could not decide which of several
candidate continuations are infinite. The OEIS entry of \oeis{A399179} asks for more terms and
suggests that the sequence might be described by a finite automaton.

\paragraph{Contributions.}
We re-implemented the algorithms of~\cite{comma}, and the decade-level version of the finite graph
of~\cite{dbts}, extended to the child graph. Our main results concern the child graph in base 10.
\begin{enumerate}[label=(\roman*)]
  \item The child graph has exactly one infinite path starting at a root; it is \oeis{A367620}
        (Theorem~\ref{thm:unique}).
  \item The branch choices of this path, \oeis{A399179}, are periodic from decade $11$ on with
        period $924$ decades, $314$ choices per period (Theorem~\ref{thm:periodic}). This gives all
        terms of \oeis{A399179}; its first $30$ terms agree with the known ones.
\end{enumerate}
We also record some smaller observations:
\begin{enumerate}[label=(\roman*),resume]
  \item The skip-ahead formulas (6.3) and (6.4) of~\cite{comma} lack a factor $b$
        (Section~\ref{sec:jumps}).
  \item Where successor sequences starting at $c \cdot 10^j$ end: the dependence on $j \bmod 924$,
        the joint distribution of start digit and end point, and the exact mean lifetime
        $(b-1)(b+2) / (2(b-2))$ decades (Section~\ref{sec:ends}).
\end{enumerate}
We also reproduce, independently, the exploration of the 50 child trees in Section~11
of~\cite{comma} (Section~\ref{sec:trees}), and the finiteness of all base-10 comma sequences
from~\cite{dbts}.

Our results rely on the periodicity statement of~\cite[Proposition~8]{dbts}, applied also to the
positions of branch points and to step counts (Section~\ref{sec:table}). Apart from that, they are
computational, and the code is described in Appendix~\ref{app:software}.
